% M11, 11.09.2026: Prüfansatz und Reichweite statt Vollreproduktions-Zusage.
% Gegenprüfung: GateResponder.Resolve; ReClarifyBacklogLoopTests;
% SubworkflowPortSpikeTests; tools/smoke-hitl.sh. Keine neuen Testläufe.
% B-30/B-32: Wiederverwendung und Replay bleiben pfadspezifisch begrenzt.

\section{Reproduzierbarkeit und Grenzen}
\label{sec:realisierung-reproduzierbarkeit-grenzen}

Die Laufartefakte aus Abschnitt~\ref{sec:realisierung-observability}
machen protokollierte Ausführungsschritte nachvollziehbar.
Automatisierte Tests ergänzen diese Beobachtung durch wiederholbare
Prüfungen ausgewählter Mechanismen. Dafür müssen Implementierungs-
und Paketstand eindeutig zugeordnet sein
(Abschnitt~\ref{sec:technischer-kontext-maf-zuordnung}). Die in
Kapitel~\ref{chap:evaluation} ausgewerteten Läufe werden zusätzlich
ihren Modellen, Gate-Betriebsarten und Eingabeversionen zugeordnet.
Auch Testergebnisse gelten für den Stand ihrer Ausführung
(Umfang wird mit dem Freeze-Stand angegeben).
% TODO: Freeze-Commit und zugehörigen Testumfang mit 6.1 abstimmen.

Die Prüfungen setzen auf mehreren Ebenen an. Einzelne
Entscheidungsregeln lassen sich ohne Modell- oder Graphaufruf testen,
etwa die Auflösung einer Liste von Gate-Elementen. Ausgewählte
MAF-Graphen werden im Testprozess mit skriptgesteuerten Ersatz-Agenten
ausgeführt. Deren vorgegebene Ergebnisse machen die Steuerung gezielt
prüfbar, beispielsweise den Übergang in eine Reparatur und das Ende
an der Versuchsschranke. Verhaltenstests halten außerdem konkrete
Frameworkeigenschaften fest, etwa die Weiterleitung einer inneren
Portanfrage aus Abschnitt~\ref{sec:realisierung-durabilitaet}.
Ändert sich das geprüfte Verhalten, wird die Abweichung im Test
sichtbar. Ein separater technischer Kurztest
(\emph{Prozessgrenzen-Smoke-Test}) prüft schließlich Start, Pause
und Wiederaufnahme in getrennten Programmaufrufen
(Abschnitt~\ref{sec:eval-resume}).

Für die Wiederholung eines Ablaufs sind außerdem seine konkrete
Graphkomposition und die Antwortregeln der beteiligten Gates
maßgeblich (Abschnitte~\ref{sec:realisierung-agent-executor-workflow}
und~\ref{sec:realisierung-human-gates}). Wiederverwendete Komponenten
oder die Betriebsart Replay allein sichern keine identischen Abläufe
und menschlichen Einzelentscheidungen.

Diese Tests belegen die tatsächlich geprüften Eigenschaften und
Zustandswirkungen unter ihren jeweiligen Bedingungen. Sie belegen
weder die fachliche Richtigkeit neu erzeugter Modellinhalte noch
eine vollständige Reproduktion des Gesamtsystems. Auch die
historischen Läufe gelten nicht als gemeinsame End-to-End-Prüfung
eines späteren Implementierungsstands.
Kapitel~\ref{chap:evaluation} führt die technischen Kontrollnachweise
mit referenzgestützten Inhaltsbewertungen und historischen
Fallprüfungen zusammen. Damit wird aus der beschriebenen
Realisierung eine Untersuchung ausgewählter Leistungen und Grenzen.
